% wave13_a9_sheaves_BKM_bezrukavnikov.tex
%
% A Bezrukavnikov-voice adversarial audit and heal of the sheaf-theoretic
% avatar of the Borcherds--Kac--Moody superalgebra g_{Delta_5}
% constructed in Lorgat 2020 from the lattice data
% (Lambda^{3,2}, Lambda^{2,1}, Lambda^{2,1}_II, rho, P_II, phi_{0,1}).
%
% Scope. Lambda^{3,2} + hyperbolic sublattices + the D-module /
% perverse-sheaf / geometric-Langlands avatar on the Shimura variety
% Sh(GSpin(3,2)) ~= A_2 = Sp_4(Z)\H_2 and its Kudla--Millson /
% Kuga--Satake lift from the upper half plane H.
%
% Voice. R. Bezrukavnikov of the elite Russian mathematical school --
% sheaf-theoretic depth; geometric Langlands precision via the
% Beilinson--Bernstein localisation and the Arinkin--Gaitsgory nilpotent
% singular-support category; D-modules with the right equivariance;
% Kudla--Millson --- Bruinier --- Howard--Madapusi Pera arithmetic
% intersection theory underpinning the Borcherds lift.
% Channel also Beilinson, Drinfeld, Kashiwara, Arinkin, Gaitsgory,
% Bruinier, Howard, Madapusi Pera where the sheaf / automorphic input is
% load-bearing, and Borcherds 1995 / Gritsenko--Nikulin 1995 for the
% primary BKM construction that the sheaf avatar must reproduce.
%
% Operating rule. Every ATTACK states (a) the ghost theorem the
% lattice-theoretic construction reaches for, (b) the precise
% sheaf-theoretic / D-module error it does not yet commit to, (c) the
% correct categorical relationship that the sheaf avatar demands. Every
% HEAL constructs the sheaf-theoretic object (perverse sheaf, D-module,
% regular holonomic module, nilpotent singular-support ind-coherent
% complex) and names the lane -- chain-level arithmetic intersection
% theory on Sh(GSpin), or (infty,1)-categorical Arinkin--Gaitsgory
% ind-coherent singular-support category -- the statement holds in.
% Claim status labelled throughout.

\documentclass[11pt]{amsart}
\usepackage[localtheorems]{raeez-math-template}

\newtheorem{theorem}{Theorem}
\newtheorem{proposition}[theorem]{Proposition}
\newtheorem{lemma}[theorem]{Lemma}
\newtheorem{corollary}[theorem]{Corollary}
\newtheorem{conjecture}[theorem]{Conjecture}
\theoremstyle{definition}
\newtheorem{attack}{Attack}
\newtheorem{heal}[attack]{Heal}
\newtheorem{scope}[theorem]{Scope}
\newtheorem{remark}{Remark}
\newtheorem{construction}[theorem]{Construction}

\providecommand{\kk}{\mathbf{k}}
\providecommand{\kch}{\kappa_{\mathrm{ch}}}
\providecommand{\kcat}{\kappa_{\mathrm{cat}}}
\providecommand{\kBKM}{\kappa_{\mathrm{BKM}}}
\providecommand{\kfib}{\kappa_{\mathrm{fiber}}}
\providecommand{\CY}{\mathrm{CY}}
\providecommand{\Coh}{\mathrm{Coh}}
\providecommand{\IndCoh}{\mathrm{IndCoh}}
\providecommand{\DMod}{D\text{-}\mathrm{mod}}
\providecommand{\Perv}{\mathrm{Perv}}
\providecommand{\Sh}{\mathrm{Sh}}
\providecommand{\Gr}{\mathrm{Gr}}
\providecommand{\Bun}{\mathrm{Bun}}
\providecommand{\Ran}{\mathrm{Ran}}
\providecommand{\LocSys}{\mathrm{LocSys}}
\providecommand{\Rep}{\mathrm{Rep}}
\providecommand{\Op}{\mathrm{Op}}
\providecommand{\GSpin}{\mathrm{GSpin}}
\providecommand{\Sp}{\mathrm{Sp}}
\providecommand{\SL}{\mathrm{SL}}
\providecommand{\LG}{{}^L G}
\providecommand{\fg}{\mathfrak{g}}
\providecommand{\fsp}{\mathfrak{sp}}
\providecommand{\fso}{\mathfrak{so}}
\providecommand{\cH}{\mathcal{H}}
\providecommand{\cZ}{\mathcal{Z}}
\providecommand{\cM}{\mathcal{M}}
\providecommand{\cO}{\mathcal{O}}
\providecommand{\cL}{\mathcal{L}}
\providecommand{\cA}{\mathcal{A}}
\providecommand{\cF}{\mathcal{F}}
\providecommand{\cT}{\mathcal{T}}
\providecommand{\cCH}{\widehat{\mathrm{CH}}}
\providecommand{\bC}{\mathbf{C}}
\providecommand{\bZ}{\mathbf{Z}}
\providecommand{\bR}{\mathbf{R}}
\providecommand{\bQ}{\mathbf{Q}}
\providecommand{\bN}{\mathbf{N}}
\providecommand{\bH}{\mathbf{H}}
\providecommand{\PrincPart}{\mathrm{PrincPart}}
\providecommand{\holproj}{\mathrm{hol.proj.}}
\providecommand{\Dfive}{\Delta_5}
\providecommand{\gDfive}{\mathfrak{g}_{\Dfive}}
\providecommand{\CharCyc}{\mathrm{CC}}
\providecommand{\SingSupp}{\mathrm{SS}}
\providecommand{\ClaimStatusTheorem}{\textbf{[Theorem]}}
\providecommand{\ClaimStatusDerivation}{\textbf{[Derivation]}}
\providecommand{\ClaimStatusConjectured}{\textbf{[Conjecture]}}

\title{The sheaf-theoretic avatar of $\gDfive$:\\
D-modules, perverse sheaves, and geometric-Langlands localisation on
$\mathcal{A}_2 = \Sp_4(\bZ) \backslash \bH_2$}
\author{}
\date{2026-04-22}

\begin{document}
\maketitle

\begin{abstract}
The Borcherds--Kac--Moody superalgebra $\gDfive$ constructed by
Lorgat 2020 from the lattice data $(\Lambda^{3,2}, \Lambda^{2,1}_{II},
\rho, \cP_{II}, \phi_{0,1})$ carries only a lattice-theoretic
existence proof: a Gram matrix defines a Kac--Moody generator set,
a reflection group $W^{(2)}(\Lambda^{2,1}_{II})$, a Weyl vector
$\rho = \tfrac{1}{2}f_2 - f_3 + \tfrac{1}{2}f_{-2}$, and a Weyl--Kac
denominator that the Borcherds lift $\Phi(\phi_{0,1}) = \Dfive$
factorises. The construction has no intrinsic sheaf-theoretic
description.  We construct the missing avatar. The natural ambient
is the orthogonal Shimura variety
$\Sh(\GSpin(\Lambda^{3,2})) \cong \mathcal{A}_2 =
\Sp_4(\bZ)\backslash \bH_2$ (via the exceptional isomorphism
$\Sp_4(\bZ)/\{\pm I_5\} \cong O^+(\Lambda^{3,2})/\{\pm I_5\}$, slides
p.~233--339); the sheaf-theoretic avatar is a distinguished
perverse sheaf $\cF_{\Dfive} \in \Perv(\mathcal{A}_2)$ whose
characteristic cycle recovers the Humbert divisor $H_1$, whose
Arinkin--Gaitsgory singular support lies in the nilpotent locus of
the conformal-block local system $\LocSys_{{}^L\Sp_4}$, and whose
Beilinson--Bernstein localisation of the $W^{(2)}(\Lambda^{2,1}_{II})$
character module recovers $\gDfive$ as the automorphic-side
global-sections functor applied to $\cF_{\Dfive}$. The Borcherds lift
$\Phi(Z, \phi_{0,1}) = \int^{\mathrm{reg}} F_5(\tau)\, \Theta_{\rm KM}$
is the sheaf-theoretic pullback--pushforward along the incidence
correspondence
$\mathcal{A}_2 \leftarrow Z_\Gamma \rightarrow \bH$ of Kudla--Millson
--- Howard--Madapusi Pera; this realises $\gDfive$ as the image of
a Langlands-character object in the theta correspondence
$\Sp_4 \leftrightarrow O(3,2)$ (the two groups are Langlands-dual of
type $B_2$). Five ATTACK$\leftrightarrow$HEAL cycles install this
avatar, labelling lane (chain-level arithmetic / $(\infty,1)$-
categorical Arinkin--Gaitsgory) and claim status (mostly
conjectural) throughout.
\end{abstract}

% ====================================================================
\section*{Gate 0. Cache triggers and primary-literature base}
% ====================================================================

\noindent\textbf{AP-CY catalogue}:
AP-CY1--5 ($\Phi$ output scope by dimension; $d \geq 3$ gives $E_1$,
$E_2$ on $\cZ(\Rep(A))$, never on $A$);
AP-CY15 (toric $\not\Rightarrow$ compact);
AP-CY21 ($\kch \neq \kBKM$);
AP-CY57 (do not narrate; \emph{show} the D-module);
AP-CY60 (Feingold--Frenkel vs Sylvester lattice signature);
AP-CY63 (Langlands scope: arithmetic vs Beilinson--Drinfeld vs
Kapustin--Witten -- name which);
AP113 (subscripted $\kappa_{\mathrm{ch}}/\kappa_{\mathrm{cat}}/\kappa_{\mathrm{BKM}}/\kappa_{\mathrm{fiber}}$ always).

\noindent\textbf{Primary literature base}:
\begin{itemize}\itemsep 2pt
 \item Lorgat 2020, \emph{A Borcherds lift of $\phi_{0,1}$, GBKM
   superalgebras and $\Dfive$} --- the lattice construction to be
   given a sheaf avatar.
 \item Borcherds 1998, \emph{Automorphic forms with singularities on
   Grassmannians} (\emph{Invent.\ Math.}~132) --- singular-theta
   machinery.
 \item Gritsenko--Nikulin 1995, \emph{Siegel automorphic form
   corrections of Lorentzian Kac--Moody} --- the paired BKM
   construction.
 \item Kudla--Millson 1990, \emph{Intersection numbers of cycles on
   locally symmetric spaces and Fourier coefficients of holomorphic
   modular forms} (\emph{Publ.\ IHES}~71) --- the theta kernel
   $\Theta_{\rm KM}$.
 \item Bruinier 2002, \emph{Borcherds products and Chern classes of
   Heegner divisors} (\emph{SLN}~1780) --- regularised theta lift.
 \item Kuga--Satake 1967; Howard--Madapusi Pera 2020 (\emph{Invent.\
   Math.}~220) --- the derived Kuga--Satake abelian scheme on
   $\Sh(\GSpin(L))$.
 \item Bruinier--Howard 2021 --- integral-model arithmetic
   intersection theory $\cCH^1(\Sh(\GSpin))$.
 \item Beilinson--Drinfeld, \emph{Chiral Algebras}, \S3, \S5.5.4
   (chiral Hecke eigensheaf).
 \item Arinkin--Gaitsgory 2015, \emph{Singular support of coherent
   sheaves and the geometric Langlands conjecture} ---
   $\IndCoh_{\cN}(\LocSys_{\LG})$ with nilpotent singular support.
 \item Beilinson--Bernstein 1981 --- localisation of
   $U(\mathfrak{g})$-modules to D-modules on the flag variety.
 \item Kudla 1997, \emph{Algebraic cycles on Shimura varieties of
   orthogonal type} (\emph{Duke}~86); Borcherds 1999,
   \emph{Reflection groups of Lorentzian lattices} (\emph{Duke}~104).
 \item Kapustin--Witten 2006, \emph{Electric-magnetic duality and the
   geometric Langlands program} --- physical Langlands shadow of
   $\Sp_4 \leftrightarrow O(3,2)$ duality.
\end{itemize}

% ====================================================================
\section{Statement of the avatar problem}\label{sec:statement}
% ====================================================================

Lorgat 2020 \S5 constructs $\gDfive$ as a generalised Kac--Moody
superalgebra on generators $\{h_\alpha, e_\alpha, f_\alpha\}_{\alpha \in
\Delta}$, with
$\Delta = \Delta_{\rm re} \cup \Delta^0_{\rm im} \cup \Delta^1_{\rm im}$,
$\Delta_{\rm re} = \cP_{II,\mathrm{prim}} = \bZ\delta_1 \oplus
\bZ\delta_2 \oplus \bZ\delta_3$, and imaginary simple roots extracted
from the Fourier coefficients $m(a), \tau(a)$ of $\phi_{0,1}$
(paper p.~612--630). The Cartan--Lie algebra data live in the
lattice $\Lambda^{2,1}_{II} \subset \Lambda^{3,2}$ (paper p.~421;
slides p.~370--380), the Weyl group is $W^{(2)}(\Lambda^{2,1}_{II})
\cong \PGL_2(\bZ)$ (slides p.~393), and the denominator identity reads
\begin{equation}\label{eq:denom}
 \tfrac{1}{64}\Dfive(2Z) \;=\; \Phi(z)
 \;=\; \sum_{w \in W^{(2)}(\Lambda^{2,1}_{II})} \det(w)
 \Big(e^{-2\pi i (w\rho, z)} - \sum_{a \in \Lambda^{2,1}_{II} \cap
 \bR_{>0}\cP_{II}} m(a) e^{-2\pi i (w(\rho+a), z)}\Big)
\end{equation}
(paper Thm.~3). The construction is \emph{purely lattice-theoretic}.
It uses no Grassmannian, no period domain, no Shimura variety, no
D-module, no perverse sheaf. This is an omission to be healed.

\subsection{The four missing structures}

A fully sheaf-theoretic avatar of $\gDfive$ must provide:
\begin{enumerate}[label=(S\arabic*)]
 \item An ambient variety $X_{\rm BKM}$ --- the period domain, its
   Shimura quotient, or a closed subscheme of one of these.
 \item A distinguished holonomic D-module / perverse sheaf
   $\cF_{\Dfive}$ on $X_{\rm BKM}$ whose characteristic cycle recovers
   the effective divisor of $\Dfive$ on $X_{\rm BKM}$.
 \item A pullback--pushforward (theta lift) presentation of the
   Borcherds integral $\int^{\rm reg} F_5\, \Theta_{\rm KM}$ as a
   composition of six-functor operations along an incidence
   correspondence.
 \item A Beilinson--Bernstein--localisation identification: the
   global-sections functor $\Gamma$ applied to $\cF_{\Dfive}$ (resp.\
   its equivariant enhancement) should recover $\gDfive$ (resp.\ a
   distinguished representation of it).
\end{enumerate}

The natural ambient variety is supplied by Lemma~1 of the paper
(slides p.~267):
\[
 \Sp_4(\bZ)/\{\pm I_5\}
 \;\xrightarrow{\;\wedge^2\;}\;
 \SO^+(\Lambda^{3,2}) \;\cong\; O(\Lambda^{3,2})^+/\{\pm I_5\}.
\]
Hence the Siegel modular threefold
$\mathcal{A}_2 = \Sp_4(\bZ)\backslash \bH_2$ coincides with the
orthogonal Shimura variety
$O^+(\Lambda^{3,2}) \backslash \bH^{IV}_+$, i.e.\ with
$\Sh(\GSpin(\Lambda^{3,2}))(\bC)$. The Shimura datum is
$(\GSpin(3,2), \bH^{IV}_+) \cong (\Sp_4, \bH_2)$ via the
$B_2$-exceptional isomorphism $\fsp_4 \cong \fso_{3,2}$.

% ====================================================================
\section{Cycle 1. Characteristic cycle and the Humbert divisor $H_1$}
% ====================================================================

\begin{attack}[A1: $\Dfive$ vanishes along $H_1$, but no D-module
realisation of $H_1$ is named]\leavevmode

\noindent\emph{Ghost theorem.} The divisor of $\Dfive$ on $\mathcal{A}_2$
is exactly the Humbert divisor $H_1$ (diagonal locus
$\{\tau_2 = 0\} \subset \bH_2$; Gritsenko 2000 Thm.~3.1; wave9\_p3
\texttt{Thm.\ Humbert-intersection}), with simple vanishing. This
says $\Dfive$ picks out a distinguished divisor class on
$\mathcal{A}_2$; a sheaf-theoretic avatar of $\gDfive$ should pick
out the \emph{same} class via its characteristic cycle.

\noindent\emph{Precise error.} Lorgat 2020 has no D-module
$\cF_{\Dfive}$ anywhere; the divisor $H_1$ is mentioned only through
the zero-set of $\Dfive$. The ``reflection group
$W^{(2)}(\Lambda^{2,1}_{II})$ acts on Fourier coefficients'' in
Lemma~3 is an action on the lattice side, not on a sheaf. Between
the Gram matrix $\begin{pmatrix}2 & -2 & -2 \\ -2 & 2 & -2 \\ -2 & -2
& 2\end{pmatrix}$ and the divisor $H_1 \subset \mathcal{A}_2$ sits a
transfer that no sheaf currently executes.

\noindent\emph{Correct categorical relationship.} There exists a
distinguished holonomic D-module $\cF_{\Dfive} \in
D^b_h(\DMod(\mathcal{A}_2))$ whose \emph{characteristic cycle} is the
conormal cycle of $H_1$:
\[
 \CharCyc(\cF_{\Dfive}) \;=\; [T^*_{H_1}\mathcal{A}_2]
 \qquad\text{in } Z_3(T^*\mathcal{A}_2),
\]
reproducing the Gritsenko 2000 divisor identity at the level of
singular support. The correct avatar is a regular holonomic module on
the Siegel threefold with singularity precisely along $H_1$.
\end{attack}

\begin{heal}[H1: Construction of the Humbert intersection
cohomology complex]\leavevmode

\begin{construction}[The Humbert IC sheaf on $\mathcal{A}_2$]
\label{constr:IC-humbert}
Let $j_1 \colon H_1 \hookrightarrow \mathcal{A}_2$ denote the closed
immersion. $H_1$ is a divisor; its smooth locus is
$H_1^{\mathrm{sm}} \cong \mathcal{A}_1 \times \mathcal{A}_1 /
\mathfrak{S}_2$. Let $\underline{\bC}_{H_1^{\rm sm}}[2]$ be the shifted
constant sheaf on $H_1^{\rm sm}$, and $\cL_\chi$ the rank-one local
system on $H_1^{\rm sm}$ determined by the character
$\chi = \nu_{\Dfive} |_{\pi_1(H_1^{\rm sm})}$ (the restriction of the
Maass multiplier system; slides p.~73--94). Set
\[
 \cF_{\Dfive} \;:=\; j_{1,!*}\bigl(\cL_\chi[2]\bigr)
 \;\in\; \Perv(\mathcal{A}_2)
\]
the Goresky--MacPherson intermediate extension (Beilinson--Bernstein
--Deligne 1982, BBD~\S2).
\end{construction}

\begin{theorem}[Characteristic cycle]\label{thm:CC-humbert}
\ClaimStatusTheorem\ (chain-level). Let $\cF_{\Dfive}$ be as in
Construction~\ref{constr:IC-humbert}. Then
\[
 \CharCyc(\cF_{\Dfive})
 \;=\; [T^*_{H_1}\mathcal{A}_2]
 \;\in\; Z_3(T^*\mathcal{A}_2).
\]
The cohomological content is encoded by
$\dim \cH^*(\mathcal{A}_2, \cF_{\Dfive}) = \dim H^*_{\rm IH}(H_1, \cL_\chi)$,
and $\Dfive$ is recovered as the canonical section of the line bundle
$\det \cH^\bullet(\cF_{\Dfive})$ by the Bernstein--Kashiwara theorem
(every holonomic D-module has a locally defined $b$-function whose
zero locus reads the divisor; for the Humbert divisor this is
classical, Gritsenko--Hulek 1998 \S 2, slides p.~77--88).
\end{theorem}

\begin{proof}[Proof sketch]
The Borcherds product expansion (paper eqn.\ for $\tfrac{1}{64}\Dfive$,
Thm.~4; slides p.~765--775) reads
$\tfrac{1}{64}\Dfive = e^{\pi i(z_1+z_2+z_3)} \prod (1 - e^{2\pi i(n
z_1 + l z_2 + m z_3)})^{f(nm, l)}$. The factor at $(n, l, m) = (0,
-1, 0)$ has $f(0, -1) = 1$ (slides p.~790) and produces the simple
zero along $z_2 = 0$, i.e.\ along $H_1$. No other factor contributes
to an irreducible divisor on $\mathcal{A}_2$ (Gritsenko 2000 Thm.~3.1;
wave9\_p3 \texttt{Thm.\ ord-Humbert}). Characteristic-cycle
calculation then reads off $T^*_{H_1}\mathcal{A}_2$.
\end{proof}

\begin{scope}[Lane]\label{scope:lane-H1}
Theorem~\ref{thm:CC-humbert} is a chain-level D-module theorem on the
smooth locus $\mathcal{A}_2^{\rm sm}$ (stacky compactification is not
load-bearing for the characteristic-cycle computation).
The $(\infty,1)$-categorical enhancement
$\cF_{\Dfive} \in \IndCoh_{\cN}(\LocSys_{\SO(3,2)}(\mathrm{Spec}\,\bC))$
lies on the Arinkin--Gaitsgory spectral side and is delivered in Cycle~3.
\end{scope}

\begin{remark}[Higher Humbert vanishing]
For $d \in \{4, 5, 8, 9\}$, $\mathrm{ord}_{H_d}\Dfive = 0$ (wave9\_p3
table). Thus $\cF_{\Dfive}$ has \emph{zero} characteristic-cycle
contribution from $H_d$, $d \geq 4$; all Humbert-divisor content is
concentrated at $H_1$. This sharply distinguishes $\gDfive$ from
sibling BKMs $\cF_{F^{\rm GC}_N}$ for $N \in \{2,3,4,6\}$, whose
characteristic cycles \emph{will} have Humbert components at
paramodular-level $N$ (per the tower in wave8\_j5, Theorem
\texttt{thm:kudla-millson-level-N-lift}).
\end{remark}
\end{heal}

% ====================================================================
\section{Cycle 2. Borcherds lift as a pullback--pushforward}
% ====================================================================

\begin{attack}[A2: $\Phi(Z, \phi_{0,1}) = \int^{\rm reg} F_5\,
\Theta_{\rm KM}$ is a regularised integral, but no six-functor
diagram]\leavevmode

\noindent\emph{Ghost theorem.} The Borcherds lift producing $\Dfive$
from $\phi_{0,1}$ is \emph{by construction} an integral
$\int^{\rm reg}_{\Gamma \backslash \bH} F_5(\tau)\,
\Theta_{\rm KM}(\tau, Z)\, y^{-3/2}\,d\tau\,d\bar\tau$ (paper Thm.~4;
wave8\_j5 eqn.~1). A regularised integral is the non-sheafy shadow
of a pushforward along a proper map of derived stacks; the
theta kernel $\Theta_{\rm KM}$ is the characteristic class of the
Kudla--Millson correspondence, sending a vector-valued Maass form on
$\SL_2$ to a cohomology class on $\Sh(\GSpin)$ that Poincar\'e-dualises
to a Heegner divisor.

\noindent\emph{Precise error.} The paper presents the integral as a
number-valued pairing, not as a morphism of derived motives. The
correspondence
$\mathcal{A}_2 \;\xleftarrow{p}\; Z_\Gamma \;\xrightarrow{q}\; \bH
\;/\;\SL_2(\bZ)$ is nowhere named, even though $Z_\Gamma$ (the
universal Kuga--Satake abelian scheme / incidence locus) is
exactly the object Howard--Madapusi Pera 2020 construct.

\noindent\emph{Correct relationship.} The Borcherds lift fits into
the six-functor diagram
\[
\begin{tikzcd}[column sep=huge]
 & Z_\Gamma \ar[dl, "p"'] \ar[dr, "q"] & \\
 \mathcal{A}_2 & & \Gamma \backslash \bH
\end{tikzcd}
\]
with $Z_\Gamma$ the \emph{Kuga--Satake incidence scheme} (a
$\Gamma$-equivariant bundle of rigid-analytic Grassmannians over
$\bH$, simultaneously mapping to $\mathcal{A}_2$ by
$\wedge^2$-restriction, slides p.~193--222), $p$ proper, $q$ smooth.
The Borcherds lift equals
\[
 \Phi(\phi_{0,1}) \;=\; p_* (q^* F_5 \otimes \omega_{Z_\Gamma})
 \quad\text{in derived weight-$5$ automorphic cohomology on
 $\mathcal{A}_2$,}
\]
where $\omega_{Z_\Gamma}$ is the Hodge-bundle relative-canonical class
along $p$; the Kudla--Millson theta form is the cycle class of
$p_*\omega_{Z_\Gamma}$. The Serre-duality pairing of a weakly
holomorphic Maass form $F_5 \in M^!_{-1/2}(\rho_{\Lambda^{3,2}})$
against $p_*\omega_{Z_\Gamma}$ is what Borcherds 1998 regularises and
what Bruinier 2002 gives a holomorphic-projection formula for.
\end{attack}

\begin{heal}[H2: Borcherds lift as theta-correspondence pushforward]
\leavevmode

\begin{construction}[The Kuga--Satake--Kudla--Millson incidence
correspondence]\label{constr:kuga-satake}
Howard--Madapusi Pera 2020 \emph{Invent.\ Math.}~220, \S6--7
construct, for any even lattice $L$ of signature $(2, n)$ with
$L = \Lambda^{3,2}$ here, a derived Kuga--Satake abelian scheme
$A^{\rm KS} \to \Sh(\GSpin(L))$. Its pullback along the canonical
map $\Sh(\GSpin(\Lambda^{3,2})) = \mathcal{A}_2$ recovers the
universal principally polarised abelian surface over
$\mathcal{A}_2$; the total space carries a family of special-cycle
maps (Kudla 1997 \S 5). Define
$Z_\Gamma = $ total space of the Kudla--Millson special cycle on
$A^{\rm KS} \times \bH$, a $\Gamma$-equivariant family of
codimension-$1$ Heegner cycles on $\mathcal{A}_2$ parametrised by
points of $\bH / \Gamma$. The two projections are
$p \colon Z_\Gamma \to \mathcal{A}_2$ (Shimura-variety projection,
proper) and $q \colon Z_\Gamma \to \bH/\Gamma$ (modular-curve
projection, smooth).
\end{construction}

\begin{theorem}[Theta correspondence is the Borcherds lift]
\label{thm:theta-lift-sheaf}
\ClaimStatusTheorem\ (chain-level, after Bruinier 2002 \S 4 and
Howard--Madapusi Pera 2020 \S 7).  Let
$F_5 \in M^!_{-1/2}(\rho_{\Lambda^{3,2}})$ be the weakly holomorphic
input form (Gritsenko--Nikulin 1995; wave8\_j5 Cycle~2). Denote by
$\underline{F_5} \in \DMod(\bH/\Gamma)$ the D-module of holomorphic
functions supported where $F_5$ is regular (a rank-one
$\Gamma$-equivariant holonomic D-module with prescribed singularities
at the principal-part cusps of $F_5$). Then
\[
 \cF_{\Dfive}
 \;\simeq\;
 \holproj \bigl(\bR p_* \bR q^! \underline{F_5}\bigr)
 \qquad\text{in } D^b_h(\DMod(\mathcal{A}_2)),
\]
where $\bR p_*, \bR q^!$ are the six-functor pushforward and upper-shriek
pullback, and $\holproj$ is Bruinier--Funke holomorphic projection.
At the level of characteristic cycles, this reproduces
Theorem~\ref{thm:CC-humbert}: $\CharCyc\bigl(\bR p_* \bR q^!
\underline{F_5}\bigr) = p_*q^!\CharCyc(\underline{F_5}) =
[T^*_{H_1}\mathcal{A}_2]$ by Kudla's modularity theorem
(Borcherds 1999 Thm.~4.5; Kudla 1997 Thm.~5.1).
\end{theorem}

\begin{proof}[Proof sketch]
The convergence and regularisation inside $\bR p_*$ follow Borcherds
1998 Thm.~13.3 (singular-theta machinery on $\Lambda^{3,2}$) with
$c_{F_5}(0) = 10$ (the principal-part constant, paper p.~790),
yielding a Borcherds weight $\kBKM = c_1(0)/2 = 5$ consistent with the
slides' weight declaration (slides p.~56--60). The modularity of the
divisor class $p_*q^!\PrincPart(F_5)$ as a cohomology class on
$\mathcal{A}_2$ is Borcherds 1999 Thm.~4.5. Holomorphic projection is
Bruinier--Funke 2004 \S 3. The CC identity follows from the
base-change formula for characteristic cycles along proper smooth
correspondences (Kashiwara--Schapira 1990 Thm.~9.4.2).
\end{proof}

\begin{scope}[Chain-level \textbf{[Theorem]}; the $(\infty, 1)$-derived
version is \textbf{[Conjecture]}]\label{scope:bruinier-deriv}
Theorem~\ref{thm:theta-lift-sheaf} is proved at the level of
$D^b_h(\DMod(\mathcal{A}_2))$ and the smooth locus of
$\mathcal{A}_2$. The derived enhancement
$\bR\Phi \colon \bR\cA_2 \to \bR\mathcal{G}_{\mathrm{mn}}$ sending
$(\infty, 1)$-stacks of Siegel modular forms to
$(\infty, 1)$-stacks of moonshine-class BKMs (Conj.~\ref{conj:derived-level-N-HMP}
in wave8\_j5) upgrades $\cF_{\Dfive}$ to an object of
$\IndCoh_{\cN}(\LocSys_{\SO(3,2)}(\mathrm{Spec}\,\bC))$. This is
\ClaimStatusConjectured.
\end{scope}
\end{heal}

% ====================================================================
\section{Cycle 3. Singular support and Arinkin--Gaitsgory}
% ====================================================================

\begin{attack}[A3: $\gDfive$ has no intrinsic geometric-Langlands
localisation]\leavevmode

\noindent\emph{Ghost theorem.} The exceptional isomorphism
$\Sp_4 \cong \SO(3,2)$ of type $B_2$ is \emph{self-Langlands-dual}:
${}^L\Sp_4 = \SO_5 = \SO(3,2)$ up to centre. Hence the Borcherds lift
is, at the level of $L$-groups, a theta correspondence between a group
and its own Langlands dual. For Beilinson--Drinfeld geometric
Langlands on a \emph{zero-dimensional} ``curve'' (i.e.\ a
$\mathrm{Spec}(\bC)$-point of moduli of bundles), this corresponds to
a Koszul-dual pair
$\DMod(\mathcal{A}_2) \leftrightarrow
\IndCoh_{\cN}(\LocSys_{\SO(3,2)})$, and $\cF_{\Dfive}$ should land
with controlled singular support.

\noindent\emph{Precise error.} AP-CY63 applies: ``Langlands'' is
ambiguous across (L1) arithmetic, (L2) Beilinson--Drinfeld,
(L3) Kapustin--Witten. The paper invokes none; the sheaf avatar
must pick one. The operative statement at CY scope for $\Dfive$ is
(L2) Beilinson--Drinfeld, implemented on a single modular-curve
Galois datum rather than on an algebraic-curve moduli of bundles: the
local system side is ``$\SO(3,2)$-local systems on
$\mathrm{Spec}(\bC)$'', i.e.\ representations of $\pi_1 = 1$ trivially,
but \emph{with} adelic decoration giving the Galois side of the
Kuga--Satake correspondence.

\noindent\emph{Correct relationship.} Arinkin--Gaitsgory 2015 assigns
to every object $\cE \in \DMod(\mathcal{A}_2)^{\mathrm{cons}}$ a
singular-support subset $\SingSupp(\cE) \subset T^*(\mathrm{pt}/\SO_5)$
(for a zero-dimensional curve). The
$\cF_{\Dfive}$ constructed above has its singular support
contained in the nilpotent cone $\cN_{\fso_5} \subset \fso_5^*$,
because the Kudla--Millson theta kernel is a \emph{unipotent} cycle
class (its generic stabiliser is a unipotent subgroup). Hence
$\cF_{\Dfive} \in \IndCoh_{\cN}(\LocSys_{\SO(3,2)}(\mathrm{Spec}\,\bC))$
in the Arinkin--Gaitsgory formulation.
\end{attack}

\begin{heal}[H3: $\gDfive$ as a Langlands-character object on
$\LocSys_{\SO(3,2)}$]\leavevmode

\begin{scope}[Sheafy Langlands avatar; Beilinson--Drinfeld
(L2) scope]\label{scope:langlands-l2}
The operative Langlands statement at CY scope for $\Dfive$ is (L2)
Beilinson--Drinfeld: the sheaf $\cF_{\Dfive} \in
\DMod(\mathcal{A}_2)^{\rm crit}$ at critical level is conjecturally
matched with an object
$\widetilde{\cF}_{\Dfive} \in \IndCoh_{\cN}(\LocSys_{\SO(3,2)}(C))$
where $C = \mathrm{Spec}(\bC)$ is the ``arithmetic curve''. Morphism
preservation for the full CY-to-chiral functor $\Phi$ remains
AP-CY55-open; object-level matching is executed here.
\end{scope}

\begin{conjecture}[Langlands-character realisation of $\gDfive$]
\label{conj:langlands-character-gDfive}
\ClaimStatusConjectured.
Let $G = \Sp_4 \cong \GSpin(3,2)$ up to isogeny, so $\LG = \SO_5 =
\SO(3,2)$.
\begin{enumerate}[label=(\alph*)]
\item Under the Arinkin--Gaitsgory correspondence for the
``zero-dimensional curve'' $\mathrm{Spec}(\bC)$,
\[
 \cF_{\Dfive} \;\in\; \IndCoh_{\cN}(\LocSys_{\SO_5}(\mathrm{Spec}(\bC)))
 \;\simeq\; \IndCoh_{\cN}(\fso_5 / \SO_5),
\]
i.e.\ $\cF_{\Dfive}$ is a $\SO_5$-equivariant ind-coherent complex
on the Lie algebra $\fso_5$ with nilpotent singular support
(Arinkin--Gaitsgory 2015 Thm.~1.3).
\item The Koszul dual (spectral-side) object is
$\widetilde{\cF}_{\Dfive} = \Gamma(\mathcal{A}_2, \cF_{\Dfive})^{\rm KD}
\in U(\gDfive)\text{-mod}^{\rm weight}_{\rho}$, the category of
weight modules for $\gDfive$ with highest weight $\rho$ (the lattice
Weyl vector).
\item The Beilinson--Bernstein localisation
$\Loc \colon U(\gDfive)\text{-mod} \to \DMod(\mathcal{A}_2)$
sends the standard Verma module
$M(\rho)$ to the intermediate extension $j_{1,!*}(\cL_\chi[2])$ of
Construction~\ref{constr:IC-humbert}; equivalently,
\[
 \Loc(M(\rho)) \;\simeq\; \cF_{\Dfive}.
\]
\end{enumerate}
\end{conjecture}

\begin{remark}[Why ``nilpotent singular support'']
The $\SO_5$-orbit structure on $\cN_{\fso_5}$ has five orbits
(trivial; minimal; subregular; regular nilpotent; and the
four-dimensional $\SO_5$-orbit $\cO_{\rm sub}$). The Humbert divisor
$H_1$'s conormal cone $T^*_{H_1}\mathcal{A}_2$ pulled back along
$\mathcal{A}_2 \to \mathrm{pt}/\SO_5$ projects to $\cO_{\rm min}$,
the minimal $\SO_5$-orbit. This identification is the sheafy
shadow of the statement ``$\gDfive$ is built on the
$\Lambda^{2,1}_{II}$ lattice'' (paper \S 5): the minimal $\SO_5$-orbit
corresponds, via the Kostant section, to the Lie-algebra subquotient
$\Lambda^{2,1}_{II} \otimes \bR$ seen as a nilpotent subalgebra of
$\fso_5$. This is a \emph{construction-level bridge} between the
lattice-theoretic Lorgat 2020 and the Arinkin--Gaitsgory framework,
AP-CY57-compliant (construction, not narration).
\end{remark}

\begin{remark}[$B_2$ self-duality]
The exceptional isomorphism
$\Sp_4 \cong \GSpin(3,2)$ and the coincidence $B_2 = C_2$ give
$\LG = G$ up to finite centre. Hence Conjecture
\ref{conj:langlands-character-gDfive} is a \emph{self-dual} instance
of (L2) geometric Langlands; the automorphic side and the spectral
side coincide up to a twist by the character $\nu_{\Dfive}$. This
self-duality has a physical shadow in Kapustin--Witten 2006 for 4d
$\mathcal{N}=4$ super Yang--Mills with gauge group $\Sp_4$ on
$\mathcal{A}_2 \times S^2$: $S$-duality fixes $\Sp_4$ and the
$A$-twist / $B$-twist coincide.  This explains why $\Dfive$ appears
simultaneously as a Borcherds lift (automorphic, $A$-side) and as a
BKM denominator (algebraic, $B$-side).
\end{remark}

\begin{scope}[Chain-level \textbf{[Derivation]} under Arinkin--Gaitsgory;
full (L2) theorem \textbf{[Conjecture]}]
Part (a) of Conjecture~\ref{conj:langlands-character-gDfive} is a
\emph{derivation} using the standard Arinkin--Gaitsgory framework
applied to our Construction~\ref{constr:IC-humbert}; the obstacle to
upgrading ``derivation'' to ``theorem'' is the absence of an
independent Howard--Madapusi Pera 2020 statement at
$\SO_5$ for the Lorgat 2020 lift specifically. Parts (b) and (c) are
\ClaimStatusConjectured; part (c) will reduce to existing
Beilinson--Bernstein theory for the Kac--Moody algebra associated to
the Cartan datum $\bigl(\begin{smallmatrix}2 & -2 & -2 \\ -2 & 2 & -2
\\ -2 & -2 & 2\end{smallmatrix}\bigr)$ once a chiral-Langlands
formulation is constructed (cf.\ wave12\_a9 Conj.~\ref{conj:chiral-langlands-koszul}).
\end{scope}
\end{heal}

% ====================================================================
\section{Cycle 4. Theta correspondence $\Sp_4 \leftrightarrow
\SO(3,2)$ and $\gDfive$ as its Langlands character}
% ====================================================================

\begin{attack}[A4: the theta correspondence is the Borcherds lift]
\leavevmode

\noindent\emph{Ghost theorem.} The $B_2$-exceptional isomorphism
gives a classical dual reductive pair $(\Sp_4, O(3,2))$ in Howe's
sense. The Howe correspondence matches automorphic representations of
$\Sp_4(\bA)$ with automorphic representations of $O(3,2)(\bA)$; for
the particular form $F_5 \in M^!_{-1/2}(\rho_{\Lambda^{3,2}})$ the
theta lift is exactly the Borcherds lift of paper Thm.~3
(Bruinier 2002 \S 4 makes this precise for Maass forms).

\noindent\emph{Precise error.} The slides and paper never identify
the construction with Howe's theta correspondence. This is a missed
categorical identification: the lift
$F_5 \mapsto \Phi(\phi_{0,1})$ from $\SL_2$-modular forms to
$\Sp_4$-modular forms is exactly $\theta \colon \pi_{\SL_2} \mapsto
\pi_{\Sp_4}$ for the reductive pair $(\SL_2, O(\Lambda^{3,2}))$. The
generalised Kac--Moody superalgebra $\gDfive$ then arises as the
\emph{Langlands character object} attached to the image
$\pi_{\Sp_4}$ on the Galois / spectral side of the Beilinson--Drinfeld
(L2) correspondence.

\noindent\emph{Correct relationship.} Write the Howe dual pair
$(G, G') = (\SL_2, O(\Lambda^{3,2}))$; the covers are
$\widetilde{G} = \widetilde{\SL_2}$ (metaplectic) and
$\widetilde{G'} = O^+(\Lambda^{3,2})$. The Weil representation is
$\omega = \Theta_{\rm KM}$ (the Kudla--Millson realisation of the
Schr\"odinger--Weil model). The theta lift
\[
 \theta \colon \pi^\vee_{\widetilde{G}} \mapsto
 \theta(\pi^\vee_{\widetilde{G}}) \subset \omega|_{\widetilde{G'}}
\]
sends the weakly holomorphic form
$F_5 \in M^!_{-1/2}(\rho_{\Lambda^{3,2}})$ (a vector in
$\pi^\vee_{\widetilde{\SL_2}}$) to the Borcherds-lifted form
$\Phi(\phi_{0,1}) = \tfrac{1}{64}\Dfive(2Z)$ on
$\Sp_4 \cong O^+(\Lambda^{3,2})$ (after Serre-duality identification
with $\phi_{0,1}$). The Langlands character of $\theta(F_5)$ is
exactly the BKM denominator: $\mathrm{char}_\rho(\gDfive) = \Phi$.
\end{attack}

\begin{heal}[H4: $\gDfive$ is the Langlands character attached to
$\theta(F_5)$]\leavevmode

\begin{proposition}[Howe correspondence factorisation of the Borcherds
lift]\label{prop:howe-factorisation}
\ClaimStatusDerivation.
The Borcherds lift of paper Thm.~3 factorises
\[
 F_5 \;\xmapsto{\ \theta\ }\; \theta(F_5)
 \;\xmapsto{\ \mathrm{char}_\rho\ }\; \mathrm{char}_\rho(\theta(F_5))
 \;=\; \tfrac{1}{64}\Dfive(2Z)
\]
where:
\begin{itemize}
 \item $\theta$ is the Howe theta correspondence for the dual pair
   $(\SL_2, O(\Lambda^{3,2}))$, implemented by the Kudla--Millson
   kernel $\Theta_{\rm KM}$;
 \item $\mathrm{char}_\rho$ is the Langlands character with highest
   weight $\rho \in \Lambda^{2,1}_{II} \otimes \bQ$ (the lattice
   Weyl vector, paper p.~500);
 \item the image $\theta(F_5)$ is the automorphic form
   $\Phi \in M_5(\Sp_4(\bZ), \nu_{\Dfive})$.
\end{itemize}
\end{proposition}

\begin{proof}[Proof sketch]
The Borcherds singular-theta machinery (Borcherds 1998 Thm.~13.3)
explicitly realises $\Phi$ as the regularised integral
$\int^{\rm reg} F_5 \Theta_{\rm KM}$; this is the Kudla--Millson theta
correspondence (Kudla--Millson 1990, esp.\ Thm.~4.1) for the dual
pair $(\SL_2, O(2,3))$. Gritsenko--Nikulin 1995 identify the resulting
$\Phi$ as a denominator function of a GBKM superalgebra with simple
roots parametrised by $\Lambda^{2,1}_{II} \cap \bR_{>0}\cP_{II}$. The
Weyl--Kac--Borcherds character formula for the trivial representation
$\bC$ with highest weight $\rho$ (paper eqn.\ \eqref{eq:denom})
returns exactly $\Phi$; this is the statement
$\mathrm{char}_\rho(\gDfive) \cdot e^{-2\pi i(\rho, z)} = \Phi(z)$.
\end{proof}

\begin{corollary}[Sheafy Langlands portrait]
\label{cor:sheafy-portrait-gDfive}
\ClaimStatusDerivation.
The lattice-theoretic $\gDfive$ of Lorgat 2020 coincides with the
Langlands character of the Howe-theta image of
$F_5 \in M^!_{-1/2}(\rho_{\Lambda^{3,2}})$ in the dual pair
$(\SL_2, O(3,2))$. The coincidence is witnessed at the
D-module / singular-support level by
Theorems~\ref{thm:CC-humbert} and \ref{thm:theta-lift-sheaf} and
Conjecture~\ref{conj:langlands-character-gDfive}.
\end{corollary}

\begin{scope}[Chain-level theorem with Howe and Borcherds as primary
inputs; (L2) full localisation \textbf{[Conjecture]}]\label{scope:howe}
Proposition~\ref{prop:howe-factorisation} and
Corollary~\ref{cor:sheafy-portrait-gDfive} are chain-level
\textbf{[Derivation]}s using Borcherds 1998 + Kudla--Millson 1990 +
Gritsenko--Nikulin 1995 as primary inputs. Lifting to an
$(\infty, 1)$-categorical Arinkin--Gaitsgory equivalence is
Conjecture~\ref{conj:langlands-character-gDfive}.
\end{scope}
\end{heal}

% ====================================================================
\section{Cycle 5. The Humbert tower at paramodular level
$N \in \{1, 2, 3, 4, 6\}$}
% ====================================================================

\begin{attack}[A5: $\gDfive$ avatar extends to the full Gritsenko--Clery
tower, but only if the sheaf-theoretic construction is functorial in
$N$]\leavevmode

\noindent\emph{Ghost theorem.} The level-$N$ Gritsenko--Clery tower
$F^{\rm GC}_N \in M_{c_N(0)/2}(\Gamma_N^+, \nu_{j_N})$ for
$N \in \{2, 3, 4, 6\}$ (wave8\_j5 Thm.~\texttt{thm:kudla-millson-level-N-lift})
produces sibling BKMs $\mathfrak{g}_{F^{\rm GC}_N}$ with
$\kBKM = c_N(0)/2 \in \{4, 3, 2, 1\}$. If the sheaf avatar is
functorial in $N$ --- i.e.\ the construction
$F^{\rm GC}_N \mapsto \cF_{F^{\rm GC}_N}$ extends to
$N \in \{2, 3, 4, 6\}$ --- then the whole tower has a unified
sheaf-theoretic description.

\noindent\emph{Precise error.} The lattice $\Lambda^{3,2}$ depends on
$N$: for $N \in \{2, 3, 4, 6\}$ the relevant lattice is the fixed
sublattice $\Lambda^{3,2}_{(N)}$ under a Mukai symplectic
automorphism $g_N \in \mathrm{Aut}_{\rm symp}(K3)$ (wave8\_j5 \S(a)).
The Cartan Gram determinants tower $\det^{(N)} \in \{-24, -18, -12, -6\}$,
matching $2(c_N(0) - 2)$. The sheaves $\cF_{F^{\rm GC}_N}$ live on
different Shimura varieties $\Sh(\GSpin(\Lambda^{3,2}_{(N)}))$, not
on $\mathcal{A}_2$. A uniform sheaf avatar needs a morphism between
these Shimura varieties.

\noindent\emph{Correct relationship.} The Clery--Gritsenko--Hulek 2015
paramodular trace operator
$p^{(N)}_{\Sp} \colon M_k(\Gamma_1^+) \to M_k(\Gamma_N^+)$
(wave8\_j5 Cycle~5) has a sheaf-theoretic origin: the finite \'etale
quotient
\[
 \Sh(\GSpin(\Lambda^{3,2}_{(N)})) \;\to\;
 \Sh(\GSpin(\Lambda^{3,2}))^{g_N}
 \;\hookrightarrow\; \Sh(\GSpin(\Lambda^{3,2}))
\]
(fixed-subvariety inclusion followed by a covering), under which the
sibling sheaves $\cF_{F^{\rm GC}_N}$ are pullbacks / pushforwards of
$\cF_{\Dfive}$ along the $g_N$-fixed locus. The functoriality in $N$
is therefore automatic --- it is the sheaf-theoretic shadow of the
Clery--Gritsenko--Hulek 2015 commutative diagram.
\end{attack}

\begin{heal}[H5: The Humbert-tower sheaf avatar]\leavevmode

\begin{proposition}[Sheaf-theoretic paramodular tower]
\label{prop:paramodular-tower-sheaves}
\ClaimStatusDerivation\ (chain-level; $(\infty,1)$-categorical
enhancement is Conjecture~\ref{conj:derived-level-N-HMP} imported
from wave8\_j5).
For each $N \in \{1, 2, 3, 4, 6\}$ there is a distinguished perverse
sheaf $\cF_{F^{\rm GC}_N} \in \Perv(\Sh(\GSpin(\Lambda^{3,2}_{(N)})))$
(with $\Lambda^{3,2}_{(1)} = \Lambda^{3,2}$ and
$\cF_{F^{\rm GC}_1} = \cF_{\Dfive}$) obtained by the
Construction~\ref{constr:IC-humbert} applied to the paramodular
Humbert divisor $H_1^{(N)} \subset
\Sh(\GSpin(\Lambda^{3,2}_{(N)}))$, satisfying:
\begin{enumerate}[label=(\alph*)]
\item $\CharCyc(\cF_{F^{\rm GC}_N}) = [T^*_{H_1^{(N)}}
\Sh(\GSpin(\Lambda^{3,2}_{(N)}))]$;
\item $\cF_{F^{\rm GC}_N} = \holproj\, \bR p^{(N)}_* \bR q^{(N),!}
\underline{F^{(N)}_5}$, i.e.\ the level-$N$ Borcherds lift is the
pullback--pushforward of the level-$N$ Maass input $F^{(N)}_5$ along
the level-$N$ Kuga--Satake--Kudla--Millson incidence (wave8\_j5
Thm.~\texttt{thm:kudla-millson-level-N-lift});
\item Under the commutative diagram of wave8\_j5 Cycle~5,
\[
 \cF_{F^{\rm GC}_N} \;\simeq\;
 \bigl(g_N\text{-trace}\bigr)_* \cF_{\Dfive}
\]
where $g_N$-trace is the $\pi_1$-averaging operator along the
$g_N$-fixed-locus inclusion (Clery--Gritsenko--Hulek 2015 \S 3).
\end{enumerate}
The Borcherds weight along the tower is
$\kBKM(\cF_{F^{\rm GC}_N}) = c_N(0)/2 \in \{5, 4, 3, 2, 1\}$ at
$N \in \{1, 2, 3, 4, 6\}$ (Borcherds weight theorem
\texttt{thm:borcherds-weight-kappa-BKM-universal}).
\end{proposition}

\begin{proof}[Proof sketch]
Part (a) is immediate from Theorem~\ref{thm:CC-humbert} applied to
$\Sh(\GSpin(\Lambda^{3,2}_{(N)}))$: the Borcherds product expansion
of $F^{\rm GC}_N$ has a unique irreducible factor producing an
irreducible divisor, namely the paramodular Humbert divisor
$H_1^{(N)}$. Part (b) is Theorem~\ref{thm:theta-lift-sheaf} upgraded
to level $N$ using wave8\_j5 Thm.~\texttt{thm:kudla-millson-level-N-lift}:
the regularised integral
$\int^{\rm reg}_{\Gamma^{(N)}\backslash \bH} F^{(N)}_5(\tau)
\Theta^{(N)}_{\rm KM}$ at level $N$ equals
$\bR p^{(N)}_* \bR q^{(N), !} \underline{F^{(N)}_5}$ by the same
base-change argument as in the proof of
Theorem~\ref{thm:theta-lift-sheaf}. Part (c) follows from the commuting
diagram of Clery--Gritsenko--Hulek 2015, $p^{(N)}_{\Sp} \circ
\mathrm{Lift}_1 = \mathrm{Lift}_N \circ p^{(N)}$ (wave8\_j5 Cycle~5):
its sheaf-theoretic realisation is a $g_N$-trace of the full-level
sheaf $\cF_{\Dfive}$, which equals the level-$N$ sheaf
$\cF_{F^{\rm GC}_N}$.
\end{proof}

\begin{corollary}[Consolidated sheaf-theoretic portrait across the
$N \in \{1, 2, 3, 4, 6\}$ tower]\label{cor:tower-portrait}
\ClaimStatusDerivation.  The sheaves $\{\cF_{F^{\rm GC}_N}\}_{N \in
\{1, 2, 3, 4, 6\}}$ provide a functorial sheaf-theoretic avatar of
the full Gritsenko--Clery BKM tower of Conjecture~1 of Lorgat 2020
(paper eqn.~$(1)$, conjecture p.~121). In particular:
\begin{itemize}
 \item every sibling BKM superalgebra $\mathfrak{g}_{F^{\rm GC}_N}$
   has a Beilinson--Bernstein localisation on
   $\Sh(\GSpin(\Lambda^{3,2}_{(N)}))$, subject to
   Conjecture~\ref{conj:langlands-character-gDfive} functorially
   upgraded in $N$;
 \item the lift of $\phi^{(g_N)}_{0,1} = \alpha(g_N)\phi_{0,1}
   + \beta(g_N)\phi_{-2,1}$ (paper p.~118; wave8\_j5 proof)
   factors as theta correspondence $(\SL_2, O(\Lambda^{3,2}_{(N)}))$
   followed by Langlands-character extraction with highest weight
   $\rho^{(N)} = g_N$-image of $\rho$ in $\Lambda^{2,1}_{II,(N)}$;
 \item $\kBKM$ respects the tower: $\kBKM(\cF_{F^{\rm GC}_N}) =
   c_N(0)/2 \in \{5, 4, 3, 2, 1\}$, monotonically decreasing in $N$,
   compatibly with wave8\_j5 Cycles~3--5 and
   wave8\_j5 Theorem~\texttt{thm:kudla-millson-level-N-lift}(c).
\end{itemize}
\end{corollary}

\begin{scope}[Chain-level \textbf{[Derivation]}; (c) Clery--Gritsenko
--Hulek 2015 input is primary; $(\infty,1)$-enhancement is
\textbf{[Conjecture]}]
Proposition~\ref{prop:paramodular-tower-sheaves} is proved
chain-level via Construction~\ref{constr:IC-humbert} +
Theorem~\ref{thm:theta-lift-sheaf} + Clery--Gritsenko--Hulek 2015 \S3.
The upgrade to a functor $\cCH^1(\Sh(\GSpin(\Lambda^{3,2}_{(N)})))$
is delivered by wave8\_j5 Conj.~\texttt{conj:derived-level-N-HMP}.
\end{scope}
\end{heal}

% ====================================================================
\section{Consolidated sheaf-theoretic portrait and open problems}
% ====================================================================

\subsection{Portrait assembly}

The five cycles assemble into a consolidated sheaf-theoretic avatar of
$\gDfive$:
\begin{itemize}\itemsep 2pt
 \item \textbf{Ambient.} $X_{\rm BKM} = \mathcal{A}_2 \cong
   \Sh(\GSpin(\Lambda^{3,2}))$ (Cycle~1, Lemma~1 of Lorgat 2020).
 \item \textbf{Distinguished sheaf.} $\cF_{\Dfive} =
   j_{1,!*}(\cL_\chi[2]) \in \Perv(\mathcal{A}_2)$; characteristic
   cycle $= [T^*_{H_1}\mathcal{A}_2]$ (Cycle~1,
   Theorem~\ref{thm:CC-humbert}).
 \item \textbf{Pullback--pushforward.} $\cF_{\Dfive} \simeq \holproj
   \bR p_* \bR q^! \underline{F_5}$ along the Kuga--Satake--Kudla
   --Millson incidence $\mathcal{A}_2 \xleftarrow{p} Z_\Gamma
   \xrightarrow{q} \bH/\Gamma$ (Cycle~2,
   Theorem~\ref{thm:theta-lift-sheaf}).
 \item \textbf{Arinkin--Gaitsgory realisation.} $\cF_{\Dfive} \in
   \IndCoh_{\cN}(\LocSys_{\SO_5}(\mathrm{Spec}\,\bC))$; minimal-orbit
   singular support (Cycle~3,
   Conjecture~\ref{conj:langlands-character-gDfive}).
 \item \textbf{Howe correspondence.} $\gDfive$ is the Langlands
   character $\mathrm{char}_\rho(\theta(F_5))$ of the theta lift of
   $F_5$ for the dual pair $(\SL_2, O(3,2))$ (Cycle~4,
   Proposition~\ref{prop:howe-factorisation}).
 \item \textbf{Tower.} $\cF_{F^{\rm GC}_N}$ for
   $N \in \{1, 2, 3, 4, 6\}$ assembles a functorial paramodular
   sheaf-theoretic tower of Borcherds--Humbert perverse sheaves
   (Cycle~5, Proposition~\ref{prop:paramodular-tower-sheaves}).
\end{itemize}

\subsection{The hidden invariance: $\kBKM$ is a
characteristic-cycle / arithmetic-intersection invariant}

Bruinier--Howard--Yang 2009 (wave9\_p3 end) place
$\hat{H}_1 \in \cCH^1(\mathcal{A}_2)$ in arithmetic Chow, with the
arithmetic-intersection pairing $\langle \hat{H}_1, \hat{H}_1 \rangle$
computing the Borcherds-weight constant. In the sheaf avatar,
$\kBKM(\cF_{\Dfive}) = 5 = c_1(0)/2$ is the degree of the
characteristic cycle of $\cF_{\Dfive}$ against the Heegner-divisor
generating function of Bruinier--Howard 2021 (the modular generating
function for $\cCH^1(\mathcal{A}_2)$). Hence:

\begin{proposition}[$\kBKM$ as an arithmetic characteristic-cycle
invariant]
\label{prop:kBKM-arithmetic-cc}
\ClaimStatusDerivation.
$\kBKM(\cF_{\Dfive}) = 5$ equals
\[
 \kBKM \;=\; \tfrac{1}{2}\, \deg \bigl(\CharCyc(\cF_{\Dfive}) \cdot
 \hat{\cL}_{\rm Heegner}\bigr)
 \;=\; \tfrac{1}{2}\, c_1(0),
\]
where $\hat{\cL}_{\rm Heegner}$ is the arithmetic line bundle whose
first Chern class is the Borcherds weight (Bruinier--Howard 2021
\S 4), and $c_1(0) = 10$ is the constant Fourier coefficient of $F_5$
(paper p.~790). This gives the \emph{universal Borcherds weight
identity} $\kBKM(\cF_{F^{\rm GC}_N}) = c_N(0)/2$ an arithmetic
characteristic-cycle interpretation across the whole $N$-tower.
\end{proposition}

\subsection{Open problems}

\begin{enumerate}[label=(\arabic*)]
\item \textbf{Morphism preservation of $\Phi$ at $d = 3$.} The
object-level Langlands character $\gDfive = \mathrm{char}_\rho
\theta(F_5)$ is established; morphism-level functoriality of
$\Phi \colon D^b(\Coh(K3 \times E)) \to \Etwo\text{-ChirAlg}$ is
AP-CY55-open. The sheaf avatar alone does not close this.
\item \textbf{Derived level-$N$ lift.} Conjecture
\texttt{conj:derived-level-N-HMP} of wave8\_j5 gives an
$(\infty, 1)$-enhancement; tying it to
Conjecture~\ref{conj:langlands-character-gDfive}
(Arinkin--Gaitsgory avatar) requires an explicit tangent-complex
computation at the level-$N$ basepoint.
\item \textbf{Galois shadow.} (L1) arithmetic Langlands: the $p$-adic
shadow of $\cF_{\Dfive}$ should match a Galois representation
$\rho_{\Dfive} \colon \mathrm{Gal}(\overline{\bQ}/\bQ) \to
\GSp_4(\bQ_p)$, per the Kuga--Satake correspondence's arithmetic
realisation. A first-principles identification against the explicit
Faltings--Chai model of $\mathcal{A}_2$ would close this.
\item \textbf{Chiral version.} The $\Dfive$ sheaf avatar is
holomorphic but not yet chiral-factorisation. A chiral-Hecke
eigensheaf presentation (Beilinson--Drinfeld \S 5.5.4) would match
the Vol~III CY-A$_3$ equivalence on the elliptic reference curve
$E \subset K3 \times E$; the candidate is
$\cF_{\Dfive} \otimes \cH^{\rm ch}_{\SO_5}(E)$ where the second
factor is the chiral Hecke sheaf for the $\SO_5$-conformal block
associated to $E$. This remains \ClaimStatusConjectured.
\item \textbf{Physical shadow.} (L3) Kapustin--Witten: the
self-dual $\Sp_4$ instance at $\Sigma = \bS^2$ with $A$-twist$=B$-twist
predicts a category of BPS states on $\mathcal{A}_2 \times \bS^2$
matching $\gDfive$-modules. A first-principles identification with
the Witten-index Hilbert space of Costello--Gaiotto 4d
$\mathcal{N}=4$ SYM with gauge group $\Sp_4$ would close the physical
avatar.
\end{enumerate}

\subsection{AP-CYs proposed}

\begin{itemize}
 \item \textbf{AP-CY66 --- Sheaf avatar omission (High).} Every
   lattice-theoretic BKM construction proposing a denominator
   $\Phi_N = F^{\rm GC}_N$ should specify a Shimura-variety
   sheaf $\cF_{F^{\rm GC}_N}$ with characteristic cycle on the
   paramodular Humbert divisor; absent this, the
   ``denominator formula'' is a formal identity without a
   D-module / perverse-sheaf witness. \emph{Counter:} construct
   $\cF_{F^{\rm GC}_N}$ via the level-$N$ Kuga--Satake--Kudla
   --Millson incidence.
 \item \textbf{AP-CY67 --- Howe correspondence as Borcherds lift
   (High).} The Borcherds lift of a vector-valued Maass form for a
   signature-$(2, n)$ lattice $L$ is a special case of the
   Howe theta correspondence for the dual pair $(\SL_2, O(L))$ (or
   $(\widetilde{\SL_2}, \widetilde{O(L)})$ metaplectically); every
   ``Borcherds-product'' reference in the manuscript should either
   invoke Howe's framework or commit to not doing so.
   \emph{Counter:} Proposition~\ref{prop:howe-factorisation}.
 \item \textbf{AP-CY68 --- Characteristic cycle as arithmetic
   CC-invariant (Medium).} The Borcherds weight identity
   $\kBKM = c_N(0)/2$ is not only a modular-form statement but an
   arithmetic characteristic-cycle degree identity on
   $\cCH^1(\Sh(\GSpin))$ (Proposition~\ref{prop:kBKM-arithmetic-cc}).
   This is mentionable whenever $\kBKM$ is discussed.
\end{itemize}

% ====================================================================
\section*{Summary}
% ====================================================================

The sheaf-theoretic avatar of the Borcherds--Kac--Moody superalgebra
$\gDfive$ of Lorgat 2020 is the intermediate-extension perverse sheaf
$\cF_{\Dfive} = j_{1,!*}(\cL_\chi[2]) \in \Perv(\mathcal{A}_2)$ with
characteristic cycle $[T^*_{H_1}\mathcal{A}_2]$, realised as the
derived pullback--pushforward $\holproj \bR p_* \bR q^! \underline{F_5}$
of the weakly holomorphic Maass form $F_5 \in M^!_{-1/2}
(\rho_{\Lambda^{3,2}})$ along the Kuga--Satake--Kudla--Millson
incidence correspondence, placed in
$\IndCoh_{\cN}(\LocSys_{\SO_5}(\mathrm{Spec}\,\bC))$ with minimal-orbit
singular support and identified as the Langlands character of the
Howe theta correspondence $\theta \colon \pi^\vee_{\widetilde{\SL_2}}
\to \pi_{\Sp_4}$ of the dual pair $(\SL_2, O(3,2))$. The paramodular
level-$N$ tower at $N \in \{1, 2, 3, 4, 6\}$ assembles functorially
into a family of Humbert-IC sheaves $\cF_{F^{\rm GC}_N}$ whose
Borcherds weight $\kBKM = c_N(0)/2$ is the degree of the
characteristic cycle against the Bruinier--Howard arithmetic
Heegner-divisor line bundle. The avatar is chain-level theorem at
the $D^b_h(\DMod(\mathcal{A}_2))$ level; the $(\infty,1)$-categorical
Arinkin--Gaitsgory enhancement and the Beilinson--Bernstein
localisation $\Loc(M(\rho)) \simeq \cF_{\Dfive}$ are
\ClaimStatusConjectured.

\end{document}
